\begin{abstract}
Severe holiday traffic congestion along Taiwan National Freeway No.~5---which connects Taipei and Yilan through the Hsuehshan Tunnel---poses persistent mobility challenges, with average daily freeway traffic volume during the 2026 Lunar New Year holiday reaching 124.0 million vehicle-kilometers (approximately $1.3\times$ the weekday annual average). Existing traffic forecasting models typically output isolated numerical speeds without regulatory context or explicit uncertainty boundaries, whereas emerging large language model (LLM) traffic assistants are prone to numerical hallucination and rarely ground responses in live roadside telemetry. Meanwhile, official mobility applications do not publish reproducible accuracy benchmarks. To bridge these gaps, this paper presents an explainable, tool-augmented congestion prediction and departure advisory agent designed under a ``prediction-as-a-tool, LLM-does-not-compute'' paradigm. Built on 5-minute electronic toll collection (ETC) gantry telemetry (TDCS M05A and \texttt{ETagPairLive.xml}) across 10 mainline segments of Freeway No.~5, the system integrates five modules: (1) domain data-aware preprocessing that filters zero-speed/zero-volume no-vehicle artifacts and enforces a verified $-3$-minute live-to-historical slot alignment; (2) histogram-based gradient boosting (\texttt{HistGradientBoosting}) models for $+15$/$+30$/$+60$-minute speed regression and isotonic-calibrated congestion risk estimation; (3) a 43-entry verified traffic-control knowledge base (high-occupancy vehicle rules and ramp closures) combined with live incident feeds; (4) a guardrailed agent enforcing deterministic numerical verification (\texttt{check\_numbers}) with safe template fallback; and (5) interactive CLI and Web interfaces. Trained on full-year 2024 data and evaluated on 2025 (365 days) and out-of-sample 2026 (311 days), the model achieves statistically significant speed MAE reductions over persistence on weekdays ($-0.39$ to $-0.75$~km/h in 2025; $-0.36$ to $-0.55$~km/h in 2026) and regular weekends ($-0.32$ to $-0.70$~km/h; $-0.31$ to $-0.50$~km/h), while gains on long holidays are smaller and year-dependent ($-0.04$~km/h, 95\% CI [$-0.17, +0.08$] at 60~min in 2026). Crucially, under a fixed-alert-budget event-level protocol (918 events in 2025; 680 in 2026), the deployed classifier does not outperform persistence for early congestion warning, demonstrating why risk scores must be distinguished from early-warning claims. Across 72 frozen agent evaluation cases (216 runs per mode), both LLM and template modes achieve 100\% deterministic guardrail compliance, with 86\% LLM adoption and zero numerical hallucinations.
\end{abstract}

\begin{IEEEkeywords}
Intelligent transportation systems, traffic congestion forecasting, explainable AI agent, tool-augmented large language models, probability calibration, Taiwan National Freeway No.~5.
\end{IEEEkeywords}
